\documentclass[10pt,a4paper]{amsart}
\usepackage{amsmath,amssymb,amsthm,graphicx}
% Retained minimal preamble for the Stacks Project chapter more-algebra.tex (More on Algebra).
% The chapter file pinned-original-source.tex carries only \input{preamble} on its first line, so no
% preamble text lives inside the chapter itself.  The Stacks Project uses one single shared preamble for
% every chapter; the snapshot stages it identically as sources/stacks-chow/preamble.tex and
% sources/stacks-etale-cohomology/preamble.tex.  The block below is copied mechanically, line for line,
% from sources/stacks-chow/preamble.tex lines 14-18 (the Xy-pic setup) and 154-217 (the native theorem
% declarations together with the mathematical macros).  Nothing was added, removed or redefined; the
% wrapper class and the amsmath/amssymb/amsthm/graphicx base packages are supplied by the assembler.
% Chapter: more-algebra.tex -- More on Algebra.
% For commutative diagrams we use Xy-pic
\usepackage[all]{xy}

% We use 2cell for 2-commutative diagrams.
\xyoption{2cell}

% Theorem environments.
%
\theoremstyle{plain}
\newtheorem{theorem}[subsection]{Theorem}
\newtheorem{proposition}[subsection]{Proposition}
\newtheorem{lemma}[subsection]{Lemma}

\theoremstyle{definition}
\newtheorem{definition}[subsection]{Definition}
\newtheorem{example}[subsection]{Example}
\newtheorem{exercise}[subsection]{Exercise}
\newtheorem{situation}[subsection]{Situation}

\theoremstyle{remark}
\newtheorem{remark}[subsection]{Remark}
\newtheorem{remarks}[subsection]{Remarks}

\numberwithin{equation}{subsection}

% Macros
%
\def\lim{\mathop{\mathrm{lim}}\nolimits}
\def\colim{\mathop{\mathrm{colim}}\nolimits}
\def\Spec{\mathop{\mathrm{Spec}}}
\def\Hom{\mathop{\mathrm{Hom}}\nolimits}
\def\Ext{\mathop{\mathrm{Ext}}\nolimits}
\def\SheafHom{\mathop{\mathcal{H}\!\mathit{om}}\nolimits}
\def\SheafExt{\mathop{\mathcal{E}\!\mathit{xt}}\nolimits}
\def\Sch{\mathit{Sch}}
\def\Mor{\mathop{\mathrm{Mor}}\nolimits}
\def\Ob{\mathop{\mathrm{Ob}}\nolimits}
\def\Sh{\mathop{\mathit{Sh}}\nolimits}
\def\NL{\mathop{N\!L}\nolimits}
\def\CH{\mathop{\mathrm{CH}}\nolimits}
\def\proetale{{pro\text{-}\acute{e}tale}}
\def\etale{{\acute{e}tale}}
\def\QCoh{\mathit{QCoh}}
\def\Ker{\mathop{\mathrm{Ker}}}
\def\Im{\mathop{\mathrm{Im}}}
\def\Coker{\mathop{\mathrm{Coker}}}
\def\Coim{\mathop{\mathrm{Coim}}}

% Boxtimes
%
\DeclareMathSymbol{\boxtimes}{\mathbin}{AMSa}{"02}

%
% Macros for moduli stacks/spaces
%
\def\QCohstack{\mathcal{QC}\!\mathit{oh}}
\def\Cohstack{\mathcal{C}\!\mathit{oh}}
\def\Spacesstack{\mathcal{S}\!\mathit{paces}}
\def\Quotfunctor{\mathrm{Quot}}
\def\Hilbfunctor{\mathrm{Hilb}}
\def\Curvesstack{\mathcal{C}\!\mathit{urves}}
\def\Polarizedstack{\mathcal{P}\!\mathit{olarized}}
\def\Complexesstack{\mathcal{C}\!\mathit{omplexes}}
% \Pic is the operator that assigns to X its picard group, usage \Pic(X)
% \Picardstack_{X/B} denotes the Picard stack of X over B
% \Picardfunctor_{X/B} denotes the Picard functor of X over B
\def\Pic{\mathop{\mathrm{Pic}}\nolimits}
\def\Picardstack{\mathcal{P}\!\mathit{ic}}
\def\Picardfunctor{\mathrm{Pic}}
\def\Deformationcategory{\mathcal{D}\!\mathit{ef}}
\usepackage{cleveref}
\crefname{originallemma}{original lemma}{original lemmas}
\Crefname{originallemma}{Original lemma}{Original lemmas}
\makeatletter
\newlabel{more-algebra-lemma-cotangent-complex-symmetric-algebra}{{07EV}{source}{}{sourcepointer-1}{}}
\newlabel{more-algebra-lemma-cotangent-complex-symmetric-algebra@cref}{{[originallemma][0][]07EV}{[1][0][]source}{}{}{}}
\makeatother
\crefname{originallemma}{original lemma}{original lemmas}
\Crefname{originallemma}{Original lemma}{Original lemmas}
\makeatletter
\newlabel{algebra-section-netherlander}{{00S0}{source}{}{sourcepointer-2}{}}
\newlabel{algebra-section-netherlander@cref}{{[originallemma][0][]00S0}{[1][0][]source}{}{}{}}
\makeatother
\crefname{originallemma}{original lemma}{original lemmas}
\Crefname{originallemma}{Original lemma}{Original lemmas}
\makeatletter
\newlabel{algebra-lemma-presentation-sym-exterior}{{00DO}{source}{}{sourcepointer-3}{}}
\newlabel{algebra-lemma-presentation-sym-exterior@cref}{{[originallemma][0][]00DO}{[1][0][]source}{}{}{}}
\makeatother
\title{The cotangent complex of a symmetric algebra presented by a free module}
\author{The Stacks Project Authors}
\date{stacks-project:more-algebra.tex:local-snapshot-20261012; GNU FDL 1.2 or later}
\begin{document}
\maketitle
\setcounter{subsection}{1}
\noindent\textit{Editorial note.} All delivered mathematical text is the author's own, copied from the pinned Stacks Project chapter file more-algebra.tex; no statement or proof step was rewritten, reordered or removed.\par\medskip
\noindent\textit{Editorial note.} The chapter file pinned-original-source.tex carries only \textbackslash{}input\{preamble\} on its first line, so no preamble text lives inside the chapter. The retained file sources/stacks-more-algebra/retained-preamble.tex is a mechanical, line-for-line copy of the project's single shared preamble as staged in the snapshot (Xy-pic setup plus theorem declarations and mathematical macros). Nothing was added, removed or redefined, and the excerpt serves only as preamble, never as a span coordinate.\par\medskip
\noindent\textit{Editorial note.} Both counted spans start at the \textbackslash{}begin line of their environment, so the statement and the author's proof are delivered whole; surrounding prose of the section is deliberately outside the counted spans and is carried as standard prerequisites instead.\par\medskip
\noindent\textit{Editorial note.} The citation fields (source\_id, source\_version, authors, source\_url, license.name, source\_title) are identical to those of every other card this lane has produced from this chapter, so the batch selector treats them as one source.\par\medskip
\noindent\textit{Allowed standard prerequisites (editorial).}\begin{itemize}
\item The naive cotangent complex of a presentation and the module of differentials of the presented algebra. the author cites Algebra, Section algebra-section-netherlander at proof line 2077. Supplies the object computed by the statement.
\item Description of the kernel of the induced map of symmetric powers as the image of the syzygy module tensored with the lower symmetric powers. the author cites Algebra, Lemma algebra-lemma-presentation-sym-exterior at proof lines 2087-2089. Gives the degreewise form of the kernel that the computation starts from.
\item Multiplicativity of the symmetric powers, so that the square of the kernel is generated by products of its graded pieces. used at proof lines 2090-2100. Allows the square of the ideal to be described degreewise.
\end{itemize}\medskip
\ifdefined\hypertarget\hypertarget{sourcepointer-1}{}\fi
\noindent\textit{Source reference.} Original lemma 07EV: focused-seeds/AGP/sources/stacks-snapshot-20261011/tags.txt, source lines 9460--9460. The lemma's own canonical tag, quoted so that readers can resolve the label inside the delivered text; it points at the counted statement itself and carries no proof content.\par\medskip
\ifdefined\hypertarget\hypertarget{sourcepointer-2}{}\fi
\noindent\textit{Source reference.} Original lemma 00S0: focused-seeds/AGP/sources/stacks-snapshot-20261011/tags.txt, source lines 1022--1022. Navigation pointer to a statement the counted proof refers to but does not contain; nothing load-bearing is delegated to it.\par\medskip
\ifdefined\hypertarget\hypertarget{sourcepointer-3}{}\fi
\noindent\textit{Source reference.} Original lemma 00DO: focused-seeds/AGP/sources/stacks-snapshot-20261011/tags.txt, source lines 504--504. Navigation pointer to a statement the counted proof refers to but does not contain; nothing load-bearing is delegated to it.\par\medskip
\begin{lemma}
\label{lemma-cotangent-complex-symmetric-algebra}
Let $A$ be a ring. Let $0 \to K \to A^{\oplus m} \to M \to 0$
be a sequence of $A$-modules. Consider the $A$-algebra
$C = \text{Sym}^*_A(M)$ with its presentation
$\alpha : A[y_1, \ldots, y_m] \to C$
coming from the surjection $A^{\oplus m} \to M$. Then
$$
\NL(\alpha) =
(K \otimes_A C \to \bigoplus\nolimits_{j = 1, \ldots, m} C \text{d}y_j)
$$
(see Algebra, Section \ref{algebra-section-netherlander})
in particular $\Omega_{C/A} = M \otimes_A C$.
\end{lemma}

\begin{proof}
Let $J = \Ker(\alpha)$. The lemma asserts that
$J/J^2 \cong K \otimes_A C$. Note that $\alpha$ is a homomorphism
of graded algebras. We will prove that in degree $d$ we have
$(J/J^2)_d = K \otimes_A C_{d - 1}$. Note that
$$
J_d = \Ker(\text{Sym}^d_A(A^{\oplus m}) \to \text{Sym}^d_A(M))
= \Im(K \otimes_A \text{Sym}^{d - 1}_A(A^{\oplus m})
\to \text{Sym}^d_A(A^{\oplus m})),
$$
see Algebra, Lemma \ref{algebra-lemma-presentation-sym-exterior}.
It follows that $(J^2)_d = \sum_{a + b = d} J_a \cdot J_b$ is the image of
$$
K \otimes_A K \otimes_A \text{Sym}^{d - 2}_A(A^{\otimes m})
\to \text{Sym}^d_A(A^{\oplus m}).
$$
The cokernel of the map $K \otimes_A \text{Sym}^{d - 2}_A(A^{\otimes m}) \to
\text{Sym}^{d - 1}_A(A^{\oplus m})$ is $\text{Sym}^{d - 1}_A(M)$ by
the lemma referenced above.
Hence it is clear that $(J/J^2)_d = J_d/(J^2)_d$ is equal to
\begin{align*}
\Coker(
K \otimes_A K \otimes_A \text{Sym}^{d - 2}_A(A^{\otimes m})
\to K \otimes_A \text{Sym}^{d - 1}_A(A^{\otimes m}))
& = K \otimes_A \text{Sym}^{d - 1}_A(M) \\
& = K \otimes_A C_{d -1}
\end{align*}
as desired.
\end{proof}

\end{document}
